\documentclass{ximera}
\title{Limit laws}
\begin{document}
\begin{abstract}
The limit laws for sums, differences, products, quotients, powers and roots, direct substitution for polynomials and rational functions, and the squeeze theorem. Textbook §2.3.
\end{abstract}
\maketitle

In a lot of cases, the best way to know the limit is to plug in the number directly.

\begin{theorem}
\label{theorem:3-5}
If $f$ is a polynomial or a rational function and $a$ is in the domain of $f$, then
\[
\lim _{x \rightarrow a} f(x)=\answer{f(a)}
\]
\end{theorem}

If we know the limit to a few functions at a point, then we can use that information to figure out the limit of their combinations!

\begin{theorem}
\label{theorem:3-6}
Suppose $c$ is some constant and suppose the following limits exist:
\[
\lim _{x \rightarrow a} f(x) \text { and } \lim _{x \rightarrow a} g(x)
\]
Then
\begin{enumerate}
  \item \[
  \lim _{x \rightarrow a}(f(x)+g(x))=\answer{\lim _{x \rightarrow a} f(x)+\lim _{x \rightarrow a} g(x)}
  \]
  \item \[
  \lim _{x \rightarrow a}(f(x)-g(x))=\answer{\lim _{x \rightarrow a} f(x)-\lim _{x \rightarrow a} g(x)}
  \]
  \item \[
  \lim _{x \rightarrow a}(c f(x))=\answer{c \lim _{x \rightarrow a} f(x)}
  \]
  \item \[
  \lim _{x \rightarrow a}(f(x) g(x))=\answer{\lim _{x \rightarrow a} f(x) \cdot \lim _{x \rightarrow a} g(x)}
  \]
  \item \[
  \lim _{x \rightarrow a}\left(\frac{f(x)}{g(x)}\right)=\answer{\frac{\lim _{x \rightarrow a} f(x)}{\lim _{x \rightarrow a} g(x)}} \text { if } \lim _{x \rightarrow a} g(x) \neq 0
  \]
  \item \[
  \lim _{x \rightarrow a}(f(x))^{n}=\answer{\left(\lim _{x \rightarrow a} f(x)\right)^{n}} \text { when } n \in \mathbb{Z}^{+}
  \]
  \item \[
  \lim _{x \rightarrow a} \sqrt[n]{f(x)}=\answer{\sqrt[n]{\lim _{x \rightarrow a} f(x)}}
  \]
\end{enumerate}
\end{theorem}

Another important rule helps us when one function is extremely similar to another function.

\begin{theorem}
\label{theorem:3-7}
If $f(x)=g(x)$ when $x \neq a$ then, if $\lim _{x \rightarrow a} f(x)$ exists, then
\[
\lim _{x \rightarrow a} g(x)=\answer{\lim _{x \rightarrow a} f(x)}
\]
\end{theorem}

% cross-ref to original Example 2.13 (in week02/limits.tex)
We saw an example of this in Example 2.13.

Two more useful properties of limits are given next.

\begin{theorem}
\label{theorem:3-8}
If $f(x) \leq g(x)$ when $x$ is near a (except possibly at $a$ ), and the limits of $f$ and $g$ both exist as $x$ approaches $a$, then
\[
\lim _{x \rightarrow a} f(x) \leq \answer{\lim _{x \rightarrow a} g(x)}
\]
\end{theorem}

\begin{theorem}[The squeeze theorem]
\label{theorem:3-9}
If $f(x) \leq g(x) \leq h(x)$ when $x$ is near a (except possibly at a) and
\[
\lim _{x \rightarrow a} f(x)=\lim _{x \rightarrow a} h(x)=L
\]
then
\[
\lim _{x t o a} g(x)=L
\]
\end{theorem}

The squeeze theorem is a super difficult theorem to use, but it is sometimes the easiest method of finding a limit!

\begin{example}
\label{example:3-10}
Let
% work: the function is determined by the text below (x^3 times cos(pi/(2x)))
\[
f(x)=\answer{x^{3} \cos \left(\frac{\pi}{2 x}\right)} .
\]
% cross-ref to original Example 2.14 (in week02/limits.tex)
Show that $\lim _{x \rightarrow 0} f(x)=0$. It would be super nice if we could just use one of our limit rules, but we found out in Example 2.14 that $\lim _{x \rightarrow 0} \cos \left(\frac{\pi}{2 x}\right)$ does not exist. So, instead we're going to have to use the squeeze theorem.

One thing to notice is that cos always ossilates between -1 and 1. In other words:
\[
-1 \leq \cos \left(\frac{\pi}{2 x}\right) \leq 1
\]
Multiply all sides by $x^{3}$ and we have
\[
\answer{-x^{3}} \leq x^{3} \cos \left(\frac{\pi}{2 x}\right) \leq \answer{x^{3}}
\]
Now we have something that looks like the squeeze theorem! Since both $x^{3}$ and $-x^{3}$ are polynomials and 0 is in their domains we have:
\[
\lim _{x \rightarrow 0} x^{3}=\answer{0} \quad \text { and } \quad \lim _{x \rightarrow 0}\left(-x^{3}\right)=\answer{0}
\]
And, by the squeeze theorem (!) we know that
\[
\lim _{x \rightarrow 0} x^{3} \cos \left(\frac{\pi}{2 x}\right)=\answer{0}
\]
\end{example}

% AUTHOR CHOICE: example chosen during conversion, change freely
\begin{exercise}
\label{exercise:3-11}
With a partner, use the squeeze theorem to show that
\[
\lim _{x \rightarrow 0} x^{2} \sin \left(\frac{1}{x}\right)=0
\]
\begin{freeResponse}
\end{freeResponse}
\end{exercise}

\end{document}
